% Build with pdfLaTeX from the repository root or this directory.
\makeatletter
\def\input@path{{./}{../../../preamble/}{book2026-2027/preamble/}}
\makeatother
\input{preamble.tex}

\geometry{letterpaper,margin=0.75in}
\hypersetup{pdftitle={Grade 10 Term 2 Outline | Discrete Probability Distributions},hidelinks}
\linespread{1.05}
\setlength{\parskip}{4pt minus 2pt}
\setlength{\emergencystretch}{2em}
\renewcommand{\thesection}{Step~\arabic{section}}

\newcolumntype{L}{>{\raggedright\arraybackslash}X}
\newcolumntype{R}[1]{>{\raggedright\arraybackslash}p{#1}}
\newcommand{\exampleheading}[1]{\par\addvspace{8pt}\textbf{#1}\par\nobreak}
\newlist{stepcriteria}{description}{1}
\setlist[stepcriteria]{leftmargin=2.8em,labelwidth=2.3em,topsep=2pt,itemsep=1pt,parsep=0pt,partopsep=0pt}

\begin{document}

\begin{center}
{\LARGE\bfseries Discrete Probability Distributions\par}
\smallskip
{\large Grade 10 \textperiodcentered{} Global Economics \textperiodcentered{} Term 2\par}
\end{center}

Model uncertain outcomes, compare return and risk, and justify an investment choice.

\section*{From uncertain returns to an investment choice}

\section{Describe uncertain returns}

Let \(X\) be Investment A's annual return, recorded as \(2\), \(4\) or \(6\) when expressed as a percentage. These possible values form its \textbf{support}. The \textbf{probability mass function (PMF)} gives the probability of each exact outcome. The \textbf{cumulative distribution function (CDF)} gives the probability that \(X\) is at or below a threshold.

\exampleheading{C1--C2 \textperiodcentered{} Read exact and cumulative probabilities}

\begin{center}
\textbf{Investment A's annual return distribution}\par\nobreak\smallskip
\begin{tabular}{@{}ccc@{}}
\toprule
\textbf{Return \(x\) (\%)} & \textbf{PMF} \(P(X=x)\) & \textbf{CDF} \(P(X\leq x)\)\\
\midrule
2 & 0.25 & 0.25\\
4 & 0.50 & 0.75\\
6 & 0.25 & 1.00\\
\bottomrule
\end{tabular}
\end{center}

All PMF entries are non-negative and \(0.25+0.50+0.25=1\). The CDF never decreases. It is \(0\) below \(2\), stays flat between outcomes and is \(1\) at and above \(6\).
\[
\begin{alignedat}{2}
&\textbf{Point}      &\qquad P(X=2)     &=0.25\\
&\textbf{Interval}   &\qquad P(2<X\leq 4) &=0.75-0.25=0.50\\
&\textbf{Cumulative} &\qquad P(X\leq 4)   &=0.75\\
&\textbf{Complement} &\qquad P(X>4)       &=1-0.75=0.25
\end{alignedat}
\]

\textbf{Criteria for this step}
\begin{stepcriteria}
\item[C1] Characterises discrete random variables in economic and financial contexts by identifying their support, probability mass function and possible outcomes.
\item[C2] Interprets PMFs and CDFs by verifying their defining properties and using them to calculate point, interval, cumulative, and complementary probabilities in financial contexts.
\end{stepcriteria}

\clearpage
\section{Count outcomes in fixed trials}

For one loan, a Bernoulli indicator is \(1\) if it defaults within a year and \(0\) if it does not. “Success” means the event we count, so here it means a default. Let \(X\) count defaults across four loans over that year.

\exampleheading{C3--C4 \textperiodcentered{} Check the binomial conditions before calculating}

\noindent
\begin{tabularx}{\linewidth}{@{}R{0.16\linewidth} L L@{}}
\toprule
\textbf{Fixed trials} & \textbf{Two outcomes} & \textbf{Model assumptions}\\
\midrule
\(n=4\) loans & Default or no default per loan & Independent outcomes, with \(p=0.10\) for every loan\\
\bottomrule
\end{tabularx}

For exactly \(k\) defaults, the binomial coefficient counts which \(k\) of the \(n\) loans default.
\[
P(X=k)=\binom{n}{k}p^k(1-p)^{n-k},
\qquad P(X=1)=4\times 0.10\times 0.90^3=0.2916.
\]

\textbf{Criteria for this step}
\begin{stepcriteria}
\item[C3] Models binary economic outcomes with Bernoulli random variables, interpreting success and failure in contexts such as default, gain, loss, or claim occurrence.
\item[C4] Applies the binomial distribution to fixed-trial financial situations by verifying its conditions and calculating probabilities using the binomial probability formula \(P(X=k)=C(n,k)\times p^k\times (1-p)^{n-k}\).
\end{stepcriteria}

\section{Count trials until success or events in an interval}

To model sales calls, assume independent trials with the same probability \(p=0.20\) of a sale on each call. Count all calls up to and including the final sale. For rare insurance claims, count events in one week. A Poisson model is reasonable if the rate is stable and events occur independently. Its parameter \(\lambda\) is the expected count in that week.

\exampleheading{C5--C7 \textperiodcentered{} What stops the count?}

{\centering
\textbf{Three models and their probabilities}\par\nobreak\smallskip
\begin{tabularx}{\linewidth}{@{}R{0.20\linewidth} L L@{}}
\toprule
\textbf{Model} & \textbf{What \(k\) counts} & \textbf{Worked example}\\
\midrule
Geometric & Calls through the first sale\newline \(k\geq 1\)
  & First sale on call 3\newline \(0.80^2\times 0.20=0.128\)\\
\addlinespace[2pt]
Negative binomial & Calls through the \(r\)th sale\newline \(k\geq r\)
  & Second sale on call 3\newline \(2\times 0.80\times 0.20^2=0.064\)\\
\addlinespace[2pt]
Poisson & Claims in one week\newline \(k=0,1,2,\ldots\)
  & No claims when \(\lambda=2\)\newline \(e^{-2}\approx 0.1353\)\\
\bottomrule
\end{tabularx}\par
}

For the second sale on call 3, exactly one of the first two calls must be a sale. Either call can be that first sale.

\textbf{Criteria for this step}
\begin{stepcriteria}
\item[C5] Models waiting-time situations with geometric distributions, calculating the probability of the first success occurring on a given trial in financial or business contexts.
\item[C6] Applies the negative binomial distribution to situations that count the number of trials needed to obtain a fixed number of successes.
\item[C7] Evaluates rare-event situations with the Poisson distribution, interpreting the parameter \(\lambda\) in contexts such as claims, arrivals, defects, or defaults within a fixed interval.
\end{stepcriteria}

\clearpage
\section{Measure return and risk, then choose a model}

\textbf{Expected return} is the long-run average over repeated years under the same distribution. \textbf{Variance} measures spread around that average in squared units. To choose a probability model for a new situation, ask what the variable counts and check the assumptions.

\exampleheading{C8--C9 \textperiodcentered{} Use Investment A's PMF, then match a model to a new count}

\[
\begin{alignedat}{2}
&\textbf{Mean} &\qquad 0.25\times 2+0.50\times 4+0.25\times 6 &=4\\
&\textbf{Variance} &\qquad 0.25(2-4)^2+0.50(4-4)^2+0.25(6-4)^2 &=2
\end{alignedat}
\]
Weight each outcome by its probability and add to get an expected return of \(4\%\). Weight the squared distances from \(4\) in the same way to get a variance of \(2\) squared percentage points.

\begin{center}
\textbf{Model selection guide}\par\nobreak\smallskip
\begin{tabularx}{\linewidth}{@{}R{0.20\linewidth} R{0.29\linewidth} L@{}}
\toprule
\textbf{Model} & \textbf{What it counts} & \textbf{What to check}\\
\midrule
Bernoulli & Success (1) or failure (0) & One trial, two outcomes, success probability \(p\)\\
\addlinespace
Binomial & Successes in \(n\) trials & Fixed \(n\), independent binary trials, constant \(p\)\\
\addlinespace
Geometric & Trials through the first success & Independent binary trials, constant \(p>0\)\\
\addlinespace
Negative binomial & Trials through the \(r\)th success & Fixed target \(r\), independent binary trials, constant \(p>0\)\\
\addlinespace
Poisson & Events in a fixed interval & Stable rate, independent event occurrences\\
\bottomrule
\end{tabularx}
\end{center}

\textbf{Choose a model.} A lender counts defaults among \(20\) loans over a year. With two outcomes per loan, independent defaults and the same default probability for every loan, use a binomial model for the number of defaults in the fixed group.

\textbf{Criteria for this step}
\begin{stepcriteria}
\item[C8] Calculates and interprets the expected value and variance of discrete probability distributions as measures of return and risk in economic and financial contexts.
\item[C9] Compares Bernoulli, binomial, negative binomial, geometric, and Poisson distributions by selecting the most appropriate model for a given financial-risk situation.
\end{stepcriteria}

\clearpage
\section{Justify an investment choice}

Screen two hypothetical investments for inclusion in a portfolio over the same one-year horizon. Compare each investment's expected return with its spread and the possibility of a loss.

\exampleheading{C10 \textperiodcentered{} A higher expected return comes with a wider spread}

\begin{center}
\textbf{Annual returns and risk for two investments}\par\nobreak\smallskip
\begin{tabular}{@{}lcc@{}}
\toprule
\textbf{Probability or measure} & \textbf{Investment A} & \textbf{Investment B}\\
\midrule
0.25 & \(2\%\) & \(-4\%\)\\
0.50 & \(4\%\) & \(6\%\)\\
0.25 & \(6\%\) & \(16\%\)\\
\midrule
Expected return & \(4\%\) & \(6\%\)\\
Variance (squared percentage points) & 2 & 50\\
\bottomrule
\end{tabular}
\end{center}

A cautious investor could include A for its smaller variance and positive return in every modelled outcome. An investor willing to accept more risk could include B for its higher expected return, accepting its larger variance and \(25\%\) chance of a loss.

\textbf{Criteria for this step}
\begin{stepcriteria}
\item[C10] Assesses investment alternatives for inclusion in a portfolio by comparing their discrete return distributions, expected returns, variances, and risk--return trade-offs, using the results to support investment portfolio decision-making.
\end{stepcriteria}

\end{document}
